%% rset.stdfmset.tex
%% A brief, self-contained sketch of Fraenkel--Mostowski set theory.
%% Drop in via %% rset.stdfmset.tex
%% A brief, self-contained sketch of Fraenkel--Mostowski set theory.
%% Drop in via %% rset.stdfmset.tex
%% A brief, self-contained sketch of Fraenkel--Mostowski set theory.
%% Drop in via %% rset.stdfmset.tex
%% A brief, self-contained sketch of Fraenkel--Mostowski set theory.
%% Drop in via \input{rset.stdfmset} (the hook already exists in rset.tex).
%%
%% NOTE ON MACROS.  The black copy of the theory is rendered here with an
%% underline (\blk) and \bin for the black membership relation, matching the
%% PDF.  If rset.local already defines your own red/black colour macros,
%% delete the \providecommand block below and the file will pick yours up.
\providecommand{\Atoms}{\mathbb{A}}        % the set of atoms
\providecommand{\Sym}{\mathrm{Sym}}        % symmetric group
\providecommand{\pow}{\mathcal{P}}         % power set
\providecommand{\Ord}{\mathrm{Ord}}        % ordinals
\providecommand{\supp}{\mathrm{supp}}      % least support
\providecommand{\fs}{\mathrm{fs}}          % finitely supported
\providecommand{\new}{\mathbin{\reflectbox{\ensuremath{\mathsf{N}}}}} % fresh quantifier

\section{Fraenkel--Mostowski set theory, in brief}
\label{sec:fm}

So that these notes stand on their own we recall the essentials of
Fraenkel--Mostowski set theory in the form made popular by Gabbay and Pitts
\cite{DBLP:journals/fac/GabbayP02}; a reader already fluent in nominal
techniques may skip directly to Section~\ref{sec:fmmodel}.

\paragraph{Atoms and the cumulative hierarchy.}
Fix a countably infinite set $\Atoms$ whose elements we call \emph{atoms}.
Atoms are \emph{not} sets: they carry no members of their own, yet they are
permitted to occur as members of sets. Over $\Atoms$ one erects the cumulative
hierarchy of $\mathsf{ZFA}$ --- Zermelo--Fraenkel set theory \emph{with atoms}
--- by transfinite recursion,
\[
  \mathcal{U}_0 = \emptyset,\qquad
  \mathcal{U}_{\alpha+1} = \Atoms \cup \pow(\mathcal{U}_\alpha),\qquad
  \mathcal{U}_\lambda = \bigcup_{\alpha<\lambda}\mathcal{U}_\alpha\ \ (\lambda\text{ a limit}),
\]
and sets $\mathcal{U} = \bigcup_{\alpha\in\Ord}\mathcal{U}_\alpha$. The axioms of
$\mathsf{ZFA}$ are those of $\mathsf{ZF}$, with extensionality, separation and
replacement relaxed to leave the membersless atoms untouched.

\paragraph{The permutation action.}
Let $\Sym(\Atoms)$ be the group of all permutations (bijections)
$\pi\colon\Atoms\to\Atoms$. Each $\pi$ acts on the entire hierarchy by
$\in$-recursion,
\[
  \pi\cdot a = \pi(a)\quad(a\in\Atoms),
  \qquad
  \pi\cdot x = \{\,\pi\cdot y \;:\; y\in x\,\}\quad(x\text{ a set}),
\]
and this is a group action: the identity acts trivially and
$(\pi\sigma)\cdot x = \pi\cdot(\sigma\cdot x)$. Intuitively $\pi$ renames atoms
uniformly, reaching arbitrarily deep into a set but never \emph{inside} an atom,
which has no interior to reach.

\paragraph{Support and finite support.}
A set $w\subseteq\Atoms$ \emph{supports} $x$ when every permutation that fixes
$w$ pointwise also fixes $x$:
\[
  \bigl(\forall a\in w.\ \pi(a)=a\bigr)\ \Longrightarrow\ \pi\cdot x = x .
\]
We call $x$ \emph{finitely supported} if some \emph{finite} $w$ supports it.
The element is \emph{equivariant} when the empty set supports it, i.e. when
$\pi\cdot x = x$ for every $\pi$; equivariant data are precisely those nameable
without reference to any particular atom.

\paragraph{The theory.}
Fraenkel--Mostowski set theory ($\mathsf{FM}$) is $\mathsf{ZFA}$ augmented by the
\emph{finite-support axiom}: \emph{every} element of the cumulative hierarchy is
hereditarily finitely supported. Equivalently, one passes to the
\emph{Fraenkel--Mostowski universe}
\[
  \mathcal{U}_{\mathsf{FM}} \;=\;
  \{\, x\in\mathcal{U} \;:\; x \text{ is hereditarily finitely supported}\,\},
\]
the largest subuniverse on which the $\Sym(\Atoms)$-action has only
finitely-supported elements. Cutting the hierarchy down in this way is exactly
the move that distinguishes $\mathsf{FM}$ from bare $\mathsf{ZFA}$.

\paragraph{Least support, freshness, and the fresh quantifier.}
Two facts give $\mathsf{FM}$ its working calculus. First, for the symmetric group
on atoms the intersection of two supports is again a support, so every
finitely-supported $x$ possesses a \emph{least} support $\supp(x)$, a canonical
finite set of atoms ``$x$ depends on''. Second, $\Atoms$ is infinite, so beyond
any finite list of parameters there is always a fresh atom. One writes
\[
  a \mathrel{\#} x \;\;:\Longleftrightarrow\;\; a\notin\supp(x)
  \qquad(\text{``$a$ is fresh for $x$''}),
\]
and introduces the \emph{fresh quantifier}
\[
  \new a.\,\phi(a) \;\;:\Longleftrightarrow\;\;
  \{\,a\in\Atoms : \neg\phi(a)\,\}\ \text{is finite},
\]
read ``for a fresh atom $a$, $\phi(a)$''. Its defining property is the
\emph{some/any theorem}: for any finitely-supported $\phi$ with parameters
$\bar p$,
\[
  \bigl(\new a.\,\phi(a)\bigr)
  \;\Longleftrightarrow\;
  \bigl(\exists a\mathrel{\#}\bar p.\ \phi(a)\bigr)
  \;\Longleftrightarrow\;
  \bigl(\forall a\mathrel{\#}\bar p.\ \phi(a)\bigr).
\]
That ``for some fresh atom'' and ``for all fresh atoms'' coincide is the
technical heart of the nominal approach to binding, and it is everything in this
list --- atoms, the permutation action, finite support, freshness --- that a
model of $\mathsf{FM}$ must reproduce. We turn to building one.
 (the hook already exists in rset.tex).
%%
%% NOTE ON MACROS.  The black copy of the theory is rendered here with an
%% underline (\blk) and \bin for the black membership relation, matching the
%% PDF.  If rset.local already defines your own red/black colour macros,
%% delete the \providecommand block below and the file will pick yours up.
\providecommand{\Atoms}{\mathbb{A}}        % the set of atoms
\providecommand{\Sym}{\mathrm{Sym}}        % symmetric group
\providecommand{\pow}{\mathcal{P}}         % power set
\providecommand{\Ord}{\mathrm{Ord}}        % ordinals
\providecommand{\supp}{\mathrm{supp}}      % least support
\providecommand{\fs}{\mathrm{fs}}          % finitely supported
\providecommand{\new}{\mathbin{\reflectbox{\ensuremath{\mathsf{N}}}}} % fresh quantifier

\section{Fraenkel--Mostowski set theory, in brief}
\label{sec:fm}

So that these notes stand on their own we recall the essentials of
Fraenkel--Mostowski set theory in the form made popular by Gabbay and Pitts
\cite{DBLP:journals/fac/GabbayP02}; a reader already fluent in nominal
techniques may skip directly to Section~\ref{sec:fmmodel}.

\paragraph{Atoms and the cumulative hierarchy.}
Fix a countably infinite set $\Atoms$ whose elements we call \emph{atoms}.
Atoms are \emph{not} sets: they carry no members of their own, yet they are
permitted to occur as members of sets. Over $\Atoms$ one erects the cumulative
hierarchy of $\mathsf{ZFA}$ --- Zermelo--Fraenkel set theory \emph{with atoms}
--- by transfinite recursion,
\[
  \mathcal{U}_0 = \emptyset,\qquad
  \mathcal{U}_{\alpha+1} = \Atoms \cup \pow(\mathcal{U}_\alpha),\qquad
  \mathcal{U}_\lambda = \bigcup_{\alpha<\lambda}\mathcal{U}_\alpha\ \ (\lambda\text{ a limit}),
\]
and sets $\mathcal{U} = \bigcup_{\alpha\in\Ord}\mathcal{U}_\alpha$. The axioms of
$\mathsf{ZFA}$ are those of $\mathsf{ZF}$, with extensionality, separation and
replacement relaxed to leave the membersless atoms untouched.

\paragraph{The permutation action.}
Let $\Sym(\Atoms)$ be the group of all permutations (bijections)
$\pi\colon\Atoms\to\Atoms$. Each $\pi$ acts on the entire hierarchy by
$\in$-recursion,
\[
  \pi\cdot a = \pi(a)\quad(a\in\Atoms),
  \qquad
  \pi\cdot x = \{\,\pi\cdot y \;:\; y\in x\,\}\quad(x\text{ a set}),
\]
and this is a group action: the identity acts trivially and
$(\pi\sigma)\cdot x = \pi\cdot(\sigma\cdot x)$. Intuitively $\pi$ renames atoms
uniformly, reaching arbitrarily deep into a set but never \emph{inside} an atom,
which has no interior to reach.

\paragraph{Support and finite support.}
A set $w\subseteq\Atoms$ \emph{supports} $x$ when every permutation that fixes
$w$ pointwise also fixes $x$:
\[
  \bigl(\forall a\in w.\ \pi(a)=a\bigr)\ \Longrightarrow\ \pi\cdot x = x .
\]
We call $x$ \emph{finitely supported} if some \emph{finite} $w$ supports it.
The element is \emph{equivariant} when the empty set supports it, i.e. when
$\pi\cdot x = x$ for every $\pi$; equivariant data are precisely those nameable
without reference to any particular atom.

\paragraph{The theory.}
Fraenkel--Mostowski set theory ($\mathsf{FM}$) is $\mathsf{ZFA}$ augmented by the
\emph{finite-support axiom}: \emph{every} element of the cumulative hierarchy is
hereditarily finitely supported. Equivalently, one passes to the
\emph{Fraenkel--Mostowski universe}
\[
  \mathcal{U}_{\mathsf{FM}} \;=\;
  \{\, x\in\mathcal{U} \;:\; x \text{ is hereditarily finitely supported}\,\},
\]
the largest subuniverse on which the $\Sym(\Atoms)$-action has only
finitely-supported elements. Cutting the hierarchy down in this way is exactly
the move that distinguishes $\mathsf{FM}$ from bare $\mathsf{ZFA}$.

\paragraph{Least support, freshness, and the fresh quantifier.}
Two facts give $\mathsf{FM}$ its working calculus. First, for the symmetric group
on atoms the intersection of two supports is again a support, so every
finitely-supported $x$ possesses a \emph{least} support $\supp(x)$, a canonical
finite set of atoms ``$x$ depends on''. Second, $\Atoms$ is infinite, so beyond
any finite list of parameters there is always a fresh atom. One writes
\[
  a \mathrel{\#} x \;\;:\Longleftrightarrow\;\; a\notin\supp(x)
  \qquad(\text{``$a$ is fresh for $x$''}),
\]
and introduces the \emph{fresh quantifier}
\[
  \new a.\,\phi(a) \;\;:\Longleftrightarrow\;\;
  \{\,a\in\Atoms : \neg\phi(a)\,\}\ \text{is finite},
\]
read ``for a fresh atom $a$, $\phi(a)$''. Its defining property is the
\emph{some/any theorem}: for any finitely-supported $\phi$ with parameters
$\bar p$,
\[
  \bigl(\new a.\,\phi(a)\bigr)
  \;\Longleftrightarrow\;
  \bigl(\exists a\mathrel{\#}\bar p.\ \phi(a)\bigr)
  \;\Longleftrightarrow\;
  \bigl(\forall a\mathrel{\#}\bar p.\ \phi(a)\bigr).
\]
That ``for some fresh atom'' and ``for all fresh atoms'' coincide is the
technical heart of the nominal approach to binding, and it is everything in this
list --- atoms, the permutation action, finite support, freshness --- that a
model of $\mathsf{FM}$ must reproduce. We turn to building one.
 (the hook already exists in rset.tex).
%%
%% NOTE ON MACROS.  The black copy of the theory is rendered here with an
%% underline (\blk) and \bin for the black membership relation, matching the
%% PDF.  If rset.local already defines your own red/black colour macros,
%% delete the \providecommand block below and the file will pick yours up.
\providecommand{\Atoms}{\mathbb{A}}        % the set of atoms
\providecommand{\Sym}{\mathrm{Sym}}        % symmetric group
\providecommand{\pow}{\mathcal{P}}         % power set
\providecommand{\Ord}{\mathrm{Ord}}        % ordinals
\providecommand{\supp}{\mathrm{supp}}      % least support
\providecommand{\fs}{\mathrm{fs}}          % finitely supported
\providecommand{\new}{\mathbin{\reflectbox{\ensuremath{\mathsf{N}}}}} % fresh quantifier

\section{Fraenkel--Mostowski set theory, in brief}
\label{sec:fm}

So that these notes stand on their own we recall the essentials of
Fraenkel--Mostowski set theory in the form made popular by Gabbay and Pitts
\cite{DBLP:journals/fac/GabbayP02}; a reader already fluent in nominal
techniques may skip directly to Section~\ref{sec:fmmodel}.

\paragraph{Atoms and the cumulative hierarchy.}
Fix a countably infinite set $\Atoms$ whose elements we call \emph{atoms}.
Atoms are \emph{not} sets: they carry no members of their own, yet they are
permitted to occur as members of sets. Over $\Atoms$ one erects the cumulative
hierarchy of $\mathsf{ZFA}$ --- Zermelo--Fraenkel set theory \emph{with atoms}
--- by transfinite recursion,
\[
  \mathcal{U}_0 = \emptyset,\qquad
  \mathcal{U}_{\alpha+1} = \Atoms \cup \pow(\mathcal{U}_\alpha),\qquad
  \mathcal{U}_\lambda = \bigcup_{\alpha<\lambda}\mathcal{U}_\alpha\ \ (\lambda\text{ a limit}),
\]
and sets $\mathcal{U} = \bigcup_{\alpha\in\Ord}\mathcal{U}_\alpha$. The axioms of
$\mathsf{ZFA}$ are those of $\mathsf{ZF}$, with extensionality, separation and
replacement relaxed to leave the membersless atoms untouched.

\paragraph{The permutation action.}
Let $\Sym(\Atoms)$ be the group of all permutations (bijections)
$\pi\colon\Atoms\to\Atoms$. Each $\pi$ acts on the entire hierarchy by
$\in$-recursion,
\[
  \pi\cdot a = \pi(a)\quad(a\in\Atoms),
  \qquad
  \pi\cdot x = \{\,\pi\cdot y \;:\; y\in x\,\}\quad(x\text{ a set}),
\]
and this is a group action: the identity acts trivially and
$(\pi\sigma)\cdot x = \pi\cdot(\sigma\cdot x)$. Intuitively $\pi$ renames atoms
uniformly, reaching arbitrarily deep into a set but never \emph{inside} an atom,
which has no interior to reach.

\paragraph{Support and finite support.}
A set $w\subseteq\Atoms$ \emph{supports} $x$ when every permutation that fixes
$w$ pointwise also fixes $x$:
\[
  \bigl(\forall a\in w.\ \pi(a)=a\bigr)\ \Longrightarrow\ \pi\cdot x = x .
\]
We call $x$ \emph{finitely supported} if some \emph{finite} $w$ supports it.
The element is \emph{equivariant} when the empty set supports it, i.e. when
$\pi\cdot x = x$ for every $\pi$; equivariant data are precisely those nameable
without reference to any particular atom.

\paragraph{The theory.}
Fraenkel--Mostowski set theory ($\mathsf{FM}$) is $\mathsf{ZFA}$ augmented by the
\emph{finite-support axiom}: \emph{every} element of the cumulative hierarchy is
hereditarily finitely supported. Equivalently, one passes to the
\emph{Fraenkel--Mostowski universe}
\[
  \mathcal{U}_{\mathsf{FM}} \;=\;
  \{\, x\in\mathcal{U} \;:\; x \text{ is hereditarily finitely supported}\,\},
\]
the largest subuniverse on which the $\Sym(\Atoms)$-action has only
finitely-supported elements. Cutting the hierarchy down in this way is exactly
the move that distinguishes $\mathsf{FM}$ from bare $\mathsf{ZFA}$.

\paragraph{Least support, freshness, and the fresh quantifier.}
Two facts give $\mathsf{FM}$ its working calculus. First, for the symmetric group
on atoms the intersection of two supports is again a support, so every
finitely-supported $x$ possesses a \emph{least} support $\supp(x)$, a canonical
finite set of atoms ``$x$ depends on''. Second, $\Atoms$ is infinite, so beyond
any finite list of parameters there is always a fresh atom. One writes
\[
  a \mathrel{\#} x \;\;:\Longleftrightarrow\;\; a\notin\supp(x)
  \qquad(\text{``$a$ is fresh for $x$''}),
\]
and introduces the \emph{fresh quantifier}
\[
  \new a.\,\phi(a) \;\;:\Longleftrightarrow\;\;
  \{\,a\in\Atoms : \neg\phi(a)\,\}\ \text{is finite},
\]
read ``for a fresh atom $a$, $\phi(a)$''. Its defining property is the
\emph{some/any theorem}: for any finitely-supported $\phi$ with parameters
$\bar p$,
\[
  \bigl(\new a.\,\phi(a)\bigr)
  \;\Longleftrightarrow\;
  \bigl(\exists a\mathrel{\#}\bar p.\ \phi(a)\bigr)
  \;\Longleftrightarrow\;
  \bigl(\forall a\mathrel{\#}\bar p.\ \phi(a)\bigr).
\]
That ``for some fresh atom'' and ``for all fresh atoms'' coincide is the
technical heart of the nominal approach to binding, and it is everything in this
list --- atoms, the permutation action, finite support, freshness --- that a
model of $\mathsf{FM}$ must reproduce. We turn to building one.
 (the hook already exists in rset.tex).
%%
%% NOTE ON MACROS.  The black copy of the theory is rendered here with an
%% underline (\blk) and \bin for the black membership relation, matching the
%% PDF.  If rset.local already defines your own red/black colour macros,
%% delete the \providecommand block below and the file will pick yours up.
\providecommand{\Atoms}{\mathbb{A}}        % the set of atoms
\providecommand{\Sym}{\mathrm{Sym}}        % symmetric group
\providecommand{\pow}{\mathcal{P}}         % power set
\providecommand{\Ord}{\mathrm{Ord}}        % ordinals
\providecommand{\supp}{\mathrm{supp}}      % least support
\providecommand{\fs}{\mathrm{fs}}          % finitely supported
\providecommand{\new}{\mathbin{\reflectbox{\ensuremath{\mathsf{N}}}}} % fresh quantifier

\section{Fraenkel--Mostowski set theory, in brief}
\label{sec:fm}

So that these notes stand on their own we recall the essentials of
Fraenkel--Mostowski set theory in the form made popular by Gabbay and Pitts
\cite{DBLP:journals/fac/GabbayP02}; a reader already fluent in nominal
techniques may skip directly to Section~\ref{sec:fmmodel}.

\paragraph{Atoms and the cumulative hierarchy.}
Fix a countably infinite set $\Atoms$ whose elements we call \emph{atoms}.
Atoms are \emph{not} sets: they carry no members of their own, yet they are
permitted to occur as members of sets. Over $\Atoms$ one erects the cumulative
hierarchy of $\mathsf{ZFA}$ --- Zermelo--Fraenkel set theory \emph{with atoms}
--- by transfinite recursion,
\[
  \mathcal{U}_0 = \emptyset,\qquad
  \mathcal{U}_{\alpha+1} = \Atoms \cup \pow(\mathcal{U}_\alpha),\qquad
  \mathcal{U}_\lambda = \bigcup_{\alpha<\lambda}\mathcal{U}_\alpha\ \ (\lambda\text{ a limit}),
\]
and sets $\mathcal{U} = \bigcup_{\alpha\in\Ord}\mathcal{U}_\alpha$. The axioms of
$\mathsf{ZFA}$ are those of $\mathsf{ZF}$, with extensionality, separation and
replacement relaxed to leave the membersless atoms untouched.

\paragraph{The permutation action.}
Let $\Sym(\Atoms)$ be the group of all permutations (bijections)
$\pi\colon\Atoms\to\Atoms$. Each $\pi$ acts on the entire hierarchy by
$\in$-recursion,
\[
  \pi\cdot a = \pi(a)\quad(a\in\Atoms),
  \qquad
  \pi\cdot x = \{\,\pi\cdot y \;:\; y\in x\,\}\quad(x\text{ a set}),
\]
and this is a group action: the identity acts trivially and
$(\pi\sigma)\cdot x = \pi\cdot(\sigma\cdot x)$. Intuitively $\pi$ renames atoms
uniformly, reaching arbitrarily deep into a set but never \emph{inside} an atom,
which has no interior to reach.

\paragraph{Support and finite support.}
A set $w\subseteq\Atoms$ \emph{supports} $x$ when every permutation that fixes
$w$ pointwise also fixes $x$:
\[
  \bigl(\forall a\in w.\ \pi(a)=a\bigr)\ \Longrightarrow\ \pi\cdot x = x .
\]
We call $x$ \emph{finitely supported} if some \emph{finite} $w$ supports it.
The element is \emph{equivariant} when the empty set supports it, i.e. when
$\pi\cdot x = x$ for every $\pi$; equivariant data are precisely those nameable
without reference to any particular atom.

\paragraph{The theory.}
Fraenkel--Mostowski set theory ($\mathsf{FM}$) is $\mathsf{ZFA}$ augmented by the
\emph{finite-support axiom}: \emph{every} element of the cumulative hierarchy is
hereditarily finitely supported. Equivalently, one passes to the
\emph{Fraenkel--Mostowski universe}
\[
  \mathcal{U}_{\mathsf{FM}} \;=\;
  \{\, x\in\mathcal{U} \;:\; x \text{ is hereditarily finitely supported}\,\},
\]
the largest subuniverse on which the $\Sym(\Atoms)$-action has only
finitely-supported elements. Cutting the hierarchy down in this way is exactly
the move that distinguishes $\mathsf{FM}$ from bare $\mathsf{ZFA}$.

\paragraph{Least support, freshness, and the fresh quantifier.}
Two facts give $\mathsf{FM}$ its working calculus. First, for the symmetric group
on atoms the intersection of two supports is again a support, so every
finitely-supported $x$ possesses a \emph{least} support $\supp(x)$, a canonical
finite set of atoms ``$x$ depends on''. Second, $\Atoms$ is infinite, so beyond
any finite list of parameters there is always a fresh atom. One writes
\[
  a \mathrel{\#} x \;\;:\Longleftrightarrow\;\; a\notin\supp(x)
  \qquad(\text{``$a$ is fresh for $x$''}),
\]
and introduces the \emph{fresh quantifier}
\[
  \new a.\,\phi(a) \;\;:\Longleftrightarrow\;\;
  \{\,a\in\Atoms : \neg\phi(a)\,\}\ \text{is finite},
\]
read ``for a fresh atom $a$, $\phi(a)$''. Its defining property is the
\emph{some/any theorem}: for any finitely-supported $\phi$ with parameters
$\bar p$,
\[
  \bigl(\new a.\,\phi(a)\bigr)
  \;\Longleftrightarrow\;
  \bigl(\exists a\mathrel{\#}\bar p.\ \phi(a)\bigr)
  \;\Longleftrightarrow\;
  \bigl(\forall a\mathrel{\#}\bar p.\ \phi(a)\bigr).
\]
That ``for some fresh atom'' and ``for all fresh atoms'' coincide is the
technical heart of the nominal approach to binding, and it is everything in this
list --- atoms, the permutation action, finite support, freshness --- that a
model of $\mathsf{FM}$ must reproduce. We turn to building one.
